\section{Experimental Results}
In this section, we present the validation results of our
methodology, followed by a comparison against two state-of-the-(SOTA)) estimation approaches. 
We also examine the overheads of synchoricity to assess the trade-offs and benefits of synchoros energy estimations.
\begin{figure}[h]
    \centering

    %--- Subfigure A ---
    \begin{subfigure}{0.48\textwidth}
        \centering
        \scriptsize
        \includesvg[scale=0.4]{figures/4_high_level_arch.svg}
        \caption{\small}
        \label{fig:arch}
    \end{subfigure}
    %\vfill
    %--- Subfigure B ---
    \begin{subfigure}{0.48\textwidth}
        \centering
        \input{figures/5_pareto}
        \caption{\small}
        \label{fig:pareto}
    \end{subfigure}

    \caption{\small (a) High-level architecture diagram with 4×4 PEs and four memory cells, and (b) pareto points for GEMM mappings, of DRRA and DiMArch CGRA ~\cite{hemani2017synchoricity,pudi2024integer,pudi2024application}.}
    \label{fig:combined}
\end{figure}
\subsection{Experimental Setup}
The energy estimation validation and evaluation are performed on DRRA and DiMArch CGRAs~\cite{hemani2017synchoricity,pudi2024integer,pudi2024application} (Fig.~\ref {fig:arch}). 
DRRA, or Dynamically Reconfigurable Resource Array, provides computational PEs consisting of a controller, a register file (RF), and an ALU for arithmetic operations.
DiMArch, or Distributed Memory Architecture, provides memory cells consisting of SRAM banks to serve as data memory. 
We choose this CGRA as the increase in computational parallelism in DRRA can be matched with an increase in parallelism in DiMArch. 
Different architectural alternatives can be implemented by clustering DRRA, DiMArch cells.
The number of DRRA columns determines the number of CGRA data memory ports.
%Note that our contributions are generally applicable for other CGRAs as well.
The cells are physically synthesized and characterized with a convergence criterion of 2\%.

%Algorithms from diverse domains are used for validation.
Table \ref{Table} lists the Berkeley Dwarf~\cite{asanovic2006landscape} and the data sizes of the algorithm workload. 
Mapping is performed by Vesyla\cite{yang2022vesyla}, the high-level synthesis tool for DRRA and DiMArch. 
These mappings are used directly to estimate the energy consumption of the resulting design.
To validate the estimates, we sweep the data sizes and all possible combinations of the number of PEs (DRRA cells) and the number of memory cells (DiMArch cells) to construct a Pareto frontier for each algorithm. 
Six points from each frontier are selected for validation (see Figure~\ref{fig:pareto}).
For ground-truth comparison, the corresponding designs are synthesized, placed, and routed in 28nm technology at 100 MHz. 
\input{tables/5_berkeley_dwarf}
\input{tables/5_parameters}
% \begin{figure}[h]
%     \centering
%     \scriptsize
%     \includesvg[scale=0.4]{figures/4_high_level_arch.svg}
% \caption{\small High-Level architecture diagram of a DRRA and DiMArch CGRA with 4x4 PEs and 4 memory cells.}
% \label{fig:arch}
% \end{figure}
% \begin{figure}[h]
%     \centering
%     \input{figures/5_pareto}
%     \caption{\small Pareto points for GEMM mappings on DRRA and DiMArch~\cite{hemani2017synchoricity,pudi2024integer,pudi2024application}.}
%     \label{fig:pareto}
% \end{figure}
\begin{figure*}[t]
    \centering
    % Subfigure (a)
    \begin{subfigure}{\textwidth}
        \centering
        \input{figures/8_power_validation}
        \caption*{\small (a)}
    \end{subfigure}
    \hfill
    % Subfigure (b)
    \begin{subfigure}{\textwidth}
        \centering
        \input{figures/9_power_validation}
        \caption*{\small (b)}
    \end{subfigure}

    \caption{\small Power estimation comparison with ground truth for three different methodologies.}
    \label{fig:power}
\end{figure*}

\begin{figure}
        \centering    
        \input{figures/6_relative}    
        \caption{\small Relative energy breakdown within a DRRA block.}        
        \label{fig:relative}
\end{figure}
\begin{figure}[t]
    \input{figures/7_granularity}
    \caption{\small Estimate accuracy at different hierarchical granularity.}
    \label{fig:granularity}
\end{figure}


% \input{tables/3_comparison}
\subsection{Validation}
We validated the total energy consumption and relative energy breakdown for each Pareto design point. 
Figure~\ref{fig:power} shows that our methodology achieves ~2\% error against ground truth across algorithm variations in functionality, dimensions, and parallelism compared against post-route ground truth results.  


Figure~\ref{fig:relative} shows the relative energy breakdown for a GEMM mapping with medium data size, 4x4 PEs, and 4 memory cells. 
The impact of all the wires is captured in the characterized energy values and, barring switching data variations, remains unchanged when performing estimations across all resources in the design.


% The validation results for variations are presented in Figure~\ref{fig:dimensions}. 
% The estimations are consistently accurate across all of the designs with average errors within 2\% relative to post-route consumption. 

\subsection{SOTA Comparison}
% We also compare our estimates with two other methodologies, Accelergy and CGRA-EAM~\cite{wijtvliet2021cgra}.
Accelergy~\cite{wu2019accelergy} specifies the design with user-defined components that are derived from basic fine-grained components. 
It focuses on datapath components, but does not consider the inter-component wires. 
For DRRA and DiMArch, the characterization database is specified in terms of 11 fine-grained components in total. 
The characterized energy values are prepared from post-layout simulations of independently synthesized fine-grained components.
Accelergy has high accuracy at the fine-grained level, such as for adders, exceeding 97\% (see Figure~\ref{fig:granularity}). 
However, accuracy decreases at higher granularity because the wires are not accounted for.
The impact of wires becomes more significant as the complexity and size of the evaluated micro-architectural units grow. 
%In contrast, our methodology is consistently accurate across all hierarchical levels in the design. 

CGRA-EAM~\cite{wijtvliet2021cgra} characterizes the power of coarse-grained components after synthesis, which includes more information details compared to the fine-grained approach. 
However, its accuracy is limited to specific characterized design instances.

\subsection{Overheads of synchoricity}
Having validated the accuracy of our methodology across functionality, dimensions, and energy breakdown, we now examine the associated overheads of synchoricity. 
In addition to DRRA+DiMArch CGRAs, we apply synchoros design principles to the HyCUBE CGRA~\cite{karunaratne2017hycube} and ADRES CGRA~\cite{mei2007adres} and quantify the design overheads introduced by synchoricity.

To apply synchoricity on designs, the PEs are synthesized and made abutment-ready after fixing the layout and peripheral interconnect locations. 
We implemented 4x4 CGRAs in 28nm process.
\input{tables/4_synchoricity}

The abutment-ready requirements in synchoricity impose restrictions on the layout of the blocks. 
Each block has to fit in the virtual grid. 
The size, shape, and location of the in and out pins are standardized to enable the composition by abutment. 
Table~\ref {tab:overhead} shows a comparison of additional overhead in terms of total standard cells and energy assuming $10\%$ switching. 
Based on the comparison, although synchoricity introduces modest overheads, we believe that the compromises are outweighed by the benefits —namely, agile and accurate energy estimation from higher levels of abstraction.
%accurate energy estimations and the reduced engineering costs of design and fabrication. 
Similar compromises were also adopted when standard cells were introduced and replaced the Mead-Conway full-custom design flow~\cite{dally201321st}.



